\rtmtight
\begin{tblr}{width=\textwidth,colspec={|Q[c,m,2.1cm]|X[l,m]|},hlines,vlines,hspan=minimal,rowsep=1.5pt,colsep=3pt,row{1}={bg=headgray,font=\bfseries}}
\SetCell{c,m}\hfil\logo\hfil & \textbf{POLITEKNIK NEGERI MALANG}\par\textbf{JURUSAN TEKNOLOGI INFORMASI}\par\textbf{PROGRAM STUDI: D4 TEKNIK INFORMATIKA} \\
\SetCell[c=2]{l} \textbf{RENCANA TUGAS MAHASISWA (RTM) DAN RUBRIK PENILAIAN} & \\
Nama Mata Kuliah & Desain Antarmuka \\
Kode Mata Kuliah & RTI252003 \quad \textbf{sks / Jam perminggu:} 2 SKS/4 jam \quad \textbf{Semester:} 2 \\
Dosen Pengampu & \begin{enumerate}\item Anugrah Nur Rahmanto, S.Sn., M.Ds.\item Farid Angga Pribadi, S.Kom., M.Kom.\item Bagas Satya Dian Nugraha, S.T., M.T.\item Dimas Wahyu Wibowo, S.T., M.T.\item Habibie Ed Dien, S.Kom., M.T.\end{enumerate} \\
\end{tblr}
\assessmentblock{Tugas Persona, Wireframe, dan Visual Design}{Case Method dan praktik studio}{SCPMK0202-01201, SCPMK0202-01202}{Mahasiswa menyelesaikan tugas persona, wireframe, dan visual design sesuai Sub-CPMK dan konteks mata kuliah.}{Menganalisis kebutuhan, mengerjakan artefak praktik/proyek, menguji hasil, mendokumentasikan evidence, dan mempresentasikan keputusan bila diminta.}{Laporan, artefak digital, repository/prototype, evidence pengujian, dan/atau bahan presentasi.}{Analisis persona dan kebutuhan pengguna & 9 & Persona, kebutuhan, pain point, dan konteks tugas dirumuskan dari data kasus serta konsisten dengan tujuan produk. & Persona dan kebutuhan utama sudah sesuai, tetapi sebagian bukti atau prioritas kebutuhan belum kuat. & Persona bersifat asumtif, tidak terkait kebutuhan pengguna, atau tidak dapat menjadi dasar desain. \\ \\
Struktur wireframe dan alur interaksi & 7.8 & Wireframe menampilkan hierarki informasi, navigasi, state utama, dan alur tugas yang runtut serta mudah diuji. & Wireframe mencakup layar utama, tetapi alur, state, atau prioritas informasi masih kurang konsisten. & Wireframe tidak menunjukkan alur tugas yang jelas atau banyak elemen tidak mendukung kebutuhan pengguna. \\ \\
Konsistensi visual design dan dokumentasi & 5.6 & Komponen visual, tipografi, warna, spacing, dan dokumentasi keputusan desain konsisten serta siap dipresentasikan. & Visual design cukup konsisten, tetapi dokumentasi keputusan atau detail komponen masih terbatas. & Visual design tidak konsisten, sulit dibaca, atau tidak terdokumentasi. \\}

\assessmentblock{UTS Review Artefak Desain Awal}{UTS berbasis presentasi artefak}{SCPMK0202-01201, SCPMK0202-01202}{Mahasiswa menyelesaikan uts review artefak desain awal sesuai Sub-CPMK dan konteks mata kuliah.}{Menganalisis kebutuhan, mengerjakan artefak praktik/proyek, menguji hasil, mendokumentasikan evidence, dan mempresentasikan keputusan bila diminta.}{Laporan, artefak digital, repository/prototype, evidence pengujian, dan/atau bahan presentasi.}{Kelengkapan artefak desain awal & 8 & Persona, user flow, wireframe, dan dasar visual tersedia lengkap serta saling terhubung. & Artefak utama tersedia, tetapi hubungan antarartefak atau detailnya belum lengkap. & Artefak desain awal tidak lengkap atau tidak cukup untuk direview. \\ \\
Kualitas evaluasi dan argumentasi desain & 7 & Keputusan desain dijelaskan dengan alasan berbasis kebutuhan pengguna, prinsip UI/UX, dan batasan kasus. & Argumentasi desain cukup jelas, tetapi sebagian keputusan belum didukung prinsip atau evidence. & Keputusan desain tidak dapat dipertahankan atau tidak terkait kebutuhan pengguna. \\ \\
Perbaikan berdasarkan umpan balik & 5 & Temuan review diprioritaskan, ditindaklanjuti, dan perubahan artefak terdokumentasi jelas. & Sebagian umpan balik ditindaklanjuti, tetapi prioritas atau dokumentasi revisi masih lemah. & Tidak ada tindak lanjut bermakna terhadap umpan balik review. \\}

\assessmentblock{Kuis Konsep dan Evaluasi UI/UX}{Kuis}{SCPMK0202-01201, SCPMK0601-01203}{Mahasiswa menyelesaikan kuis konsep dan evaluasi ui/ux sesuai Sub-CPMK dan konteks mata kuliah.}{Menganalisis kebutuhan, mengerjakan artefak praktik/proyek, menguji hasil, mendokumentasikan evidence, dan mempresentasikan keputusan bila diminta.}{Laporan, artefak digital, repository/prototype, evidence pengujian, dan/atau bahan presentasi.}{Penguasaan konsep UI/UX & 4 & Jawaban menunjukkan pemahaman prinsip usability, affordance, feedback, consistency, dan accessibility pada kasus. & Sebagian konsep dijelaskan benar, tetapi penerapannya pada kasus belum lengkap. & Jawaban keliru atau hanya menghafal istilah tanpa penerapan. \\ \\
Ketepatan evaluasi antarmuka & 3.5 & Masalah UI/UX diidentifikasi tepat, diprioritaskan, dan dikaitkan dengan dampak pada pengguna. & Identifikasi masalah cukup tepat, tetapi prioritas atau dampaknya belum jelas. & Masalah yang dipilih tidak relevan atau tidak menunjukkan evaluasi antarmuka. \\ \\
Rekomendasi perbaikan desain & 2.5 & Rekomendasi spesifik, dapat diterapkan, dan konsisten dengan prinsip UI/UX serta konteks pengguna. & Rekomendasi cukup relevan, tetapi masih umum atau belum lengkap. & Rekomendasi tidak operasional atau tidak memperbaiki masalah utama. \\}

\assessmentblock{PBL Proyek UI/UX Terintegrasi}{Project Based Learning}{SCPMK0202-01201-SCPMK0802-01204}{Mahasiswa menyelesaikan pbl proyek ui/ux terintegrasi sesuai Sub-CPMK dan konteks mata kuliah.}{Menganalisis kebutuhan, mengerjakan artefak praktik/proyek, menguji hasil, mendokumentasikan evidence, dan mempresentasikan keputusan bila diminta.}{Laporan, artefak digital, repository/prototype, evidence pengujian, dan/atau bahan presentasi.}{Kesesuaian solusi dengan riset pengguna & 9 & Solusi desain menunjukkan keterkaitan kuat antara riset, persona, journey, kebutuhan, dan prioritas fitur. & Keterkaitan riset dan solusi cukup terlihat, tetapi beberapa fitur belum didukung evidence. & Solusi tidak berangkat dari riset pengguna atau kebutuhan utama tidak terjawab. \\ \\
Kualitas prototype dan design system & 7.9 & Prototype interaktif, alur utama lengkap, komponen konsisten, dan design system dapat dipakai ulang. & Prototype berjalan untuk alur utama, tetapi konsistensi komponen atau state masih terbatas. & Prototype tidak menunjukkan alur utama atau komponen desain tidak konsisten. \\ \\
Validasi usability dan presentasi keputusan & 5.7 & Uji usability terdokumentasi, temuan dianalisis, revisi diprioritaskan, dan keputusan dipresentasikan meyakinkan. & Validasi dilakukan, tetapi analisis temuan atau tindak lanjut revisi belum kuat. & Tidak ada validasi pengguna yang memadai atau presentasi tidak menjelaskan keputusan desain. \\}

\assessmentblock{UAS Validasi Akhir Prototype}{UAS praktik dan defense}{SCPMK0601-01203, SCPMK0802-01204}{Mahasiswa menyelesaikan uas validasi akhir prototype sesuai Sub-CPMK dan konteks mata kuliah.}{Menganalisis kebutuhan, mengerjakan artefak praktik/proyek, menguji hasil, mendokumentasikan evidence, dan mempresentasikan keputusan bila diminta.}{Laporan, artefak digital, repository/prototype, evidence pengujian, dan/atau bahan presentasi.}{Kelengkapan validasi prototype & 10 & Skenario uji, partisipan, metrik, hasil observasi, dan bukti prototype akhir terdokumentasi lengkap. & Validasi tersedia, tetapi metrik, bukti observasi, atau cakupan skenario belum lengkap. & Validasi tidak terstruktur atau tidak membuktikan kualitas prototype. \\ \\
Analisis temuan dan revisi desain & 8.8 & Temuan usability dianalisis menjadi prioritas revisi yang jelas dan terlihat pada prototype akhir. & Sebagian temuan ditindaklanjuti, tetapi prioritas atau alasan revisi belum kuat. & Temuan tidak dianalisis atau revisi tidak terkait hasil validasi. \\ \\
Defense keputusan desain akhir & 6.2 & Mahasiswa mampu mempertahankan keputusan desain akhir dengan evidence, prinsip UI/UX, dan batasan implementasi. & Defense cukup jelas, tetapi beberapa keputusan belum didukung evidence. & Tidak mampu menjelaskan alasan desain atau bukti validasi akhir. \\}

\begin{tblr}{width=\textwidth,colspec={|Q[l,m,3.2cm]|X[l,m]|},hlines,vlines,hspan=minimal,row{1-3}={bg=headgray},rowsep=1.5pt,colsep=3pt}
Jadwal Pelaksanaan: & Minggu 1--17 sesuai RPS. Setiap evaluasi memiliki RTM dan rubrik multi-kriteria. \\
Lain-Lain Yang Diperlukan: & LMS, repositori/prototype/proyek, perangkat lunak sesuai mata kuliah, dan perangkat presentasi. \\
Pustaka: & Mengacu pustaka utama dan pendukung pada bagian RPS. \\
\end{tblr}
